\chapter{RESUMO}\label{chap:resumo}

A crescente adoção de dispositivos de Internet das Coisas (IoT) tem ampliado a superfície de ataque das redes atuais, uma vez que esses dispositivos costumam apresentar recursos computacionais limitados e mecanismos de segurança pouco robustos, tornando-se alvos recorrentes de ataques como varreduras de portas, negação de serviço distribuída (DDoS) e infecções por botnets como Mirai. Diante da dificuldade de abordagens tradicionais de detecção de intrusão em acompanhar a diversidade desse tráfego, este trabalho tem como objetivo comparar o desempenho de diferentes algoritmos de aprendizado de máquina na detecção de ataques em redes IoT, utilizando o conjunto de dados IoT-23, amplamente empregado na área por reunir capturas reais de tráfego benigno e malicioso. A revisão de literatura mapeou a arquitetura típica de redes IoT, suas principais vulnerabilidades e os algoritmos mais utilizados em trabalhos correlatos — Random Forest, Support Vector Machine (SVM), Árvore de Decisão e Redes Neurais Convolucionais (CNN) —, evidenciando uma lacuna relacionada à escassez de comparações sistemáticas entre esses algoritmos sob as mesmas condições experimentais sobre o IoT-23. Para preencher essa lacuna, a metodologia proposta prevê a seleção de cenários representativos do dataset, pré-processamento e balanceamento das classes, seguidos do treinamento e avaliação dos quatro algoritmos em tarefas de classificação binária e multiclasse, com desempenho comparado por meio de métricas como acurácia, precisão, recall, F1-Score e tempo de inferência, complementadas por teste de significância estatística. Espera-se, com os resultados, contribuir para a identificação do algoritmo mais adequado à detecção de ataques em redes IoT, considerando tanto a eficácia da classificação quanto a viabilidade de aplicação em ambientes com recursos computacionais restritos.